\section{Background} \label{sec:background}

\subsection{Tricotyledon Theory of System Design}
Wymore's T3SD organizes system design into three cotyledons (design spaces): Functionality (FC), Buildability (BC), and Implementability (IC) \parencite{Wymore1993}. The FC contains system designs that satisfy required input/output behavior as well as performance orderings. BC captures technology constraints and cost orderings; IC is the space of designs that homomorphically map buildable artifacts to functional ones. In classical verification, functional intent is realized constructively: system requirements are formalized into a reference system $Z_{\text{spec}}$ (a specification-system in the FC), and implementations are checked via Wymore system homomorphisms. This paper focuses specifically on the FC.

Figure~\ref{fig:t3sd-cotyledons} situates both verification readings within the cotyledon layout: $Z_{\text{spec}}$ and property set $\Phi$ express functional intent in the FC. Under classical T3SD, a verified implementation lies in the implementability overlap (i.e., buildable in BC and homomorphic to $Z_{\text{spec}}$), with $Z_{\text{impl}}$ shown in the IC region. Assertional checking instead asks whether $Z_{\text{impl}} \models_{\mathcal{L}} \Phi$ and does not require a homomorphism.

\begin{figure}[t]
  \centering
  \resizebox{0.6\textwidth}{!}{%
  \begin{tikzpicture}[
    cotyledon/.style={draw, thick, circle, minimum size=6.2cm, fill opacity=0.05},
    box/.style={draw, thick, rounded corners=4pt, align=center, fill=white, fill opacity=1, text opacity=1, minimum width=2.6cm, minimum height=0.95cm, font=\small},
    cotarrow/.style={-{Stealth[scale=1.2]}, thick},
    classic/.style={dashed, draw=black!55, text=black!75},
    modern/.style={draw=blue!80!black, thick, fill=blue!5},
    labelStyle/.style={font=\normalsize\bfseries, align=center},
    labelLine/.style={dashed, thick, draw=black!45}
  ]
    \coordinate (ICoverlap) at (2.5, {sqrt(9.61-6.25)});
    \node[cotyledon, fill=blue, draw=blue!70!black] (FC) at (0,0) {};
    \node[cotyledon, fill=orange, draw=orange!70!black] (BC) at (5,0) {};

    \node[labelStyle, text=blue!80!black, anchor=north] (FClabel) at (-1.25, -3.35) {Functionality\\Cotyledon (FC)};
    \node[labelStyle, text=orange!80!black, anchor=north] (BClabel) at (6.25, -3.35) {Buildability\\Cotyledon};
    \node[labelStyle, text=black!70, anchor=south] (IClabel) at (2.5, 3.35) {Implementability\\Cotyledon};
    \draw[labelLine] (IClabel.south) -- (ICoverlap);

    \node[box, classic, font=\scriptsize] (Zspec) at (-1.35, 0.85) {Specification\\System $Z_{\text{spec}}$};
    \node[box, modern, font=\scriptsize] (Phi) at (-1.35, -1.35) {Property Set $\Phi$\\[-1pt] {\tiny FO-LTL, STL, \ldots}};

    \node[box, draw=orange!80!black, fill=orange!5] (Zimpl) at (2.95, 0.15) {Implementation\\System $Z_{\text{impl}}$};

    \node[box, draw=orange!60!black, fill=orange!3, dashed, font=\scriptsize, minimum width=2.1cm, minimum height=0.75cm] (Zprime) at (5.55, -0.95) {$Z' \in \text{BC}$\\[-2pt]{\tiny no $h$ to $Z_{\text{spec}}$}};

    % Classical: implementation verifies via homomorphism to Z_spec
    \draw[-{Stealth[scale=0.9]}, semithick, classic]
      (Zimpl.west) to[out=180, in=0, looseness=1.05] ($(Zspec.east)+(0,0.1)$)
      node[midway, above, sloped, font=\footnotesize, text=black!75] {classical};

    % Classical spec may abstract to the same functional intent as Phi
    \draw[-{Stealth[scale=0.85]}, semithick, classic, densely dashed]
      (Zspec.south) -- (Phi.north);

    % Assertional: implementation checked against Phi (enter east face, tangent-aligned)
    \draw[-{Stealth[scale=0.9]}, semithick, modern]
      (Zimpl.west) to[out=180, in=0, looseness=1.05] ($(Phi.east)+(0,0.2)$)
      node[midway, below, sloped, font=\footnotesize, text=blue!90!black] {proposed};

    \draw[-{Stealth[scale=0.9]}, semithick, modern, dashed]
      (Zprime.west) to[out=185, in=-12, looseness=1.1] ($(Phi.east)+(0,-0.22)$);

    \node[font=\tiny, text=black!75, align=center] at (-1.35, 1.55) {Classical FC:\\[-1pt] $\exists h: Z_{\text{impl}} \to Z_{\text{spec}}$};
    \node[font=\tiny, text=blue!90!black, align=center] at (-1.35, -2.05) {Assertional FC:\\[-1pt] $\mathcal{FC}_{\mathcal{L}}(\Phi)$};
  \end{tikzpicture}%
  }
  \caption{T3SD cotyledon layout (design-space placement). Functional intent is expressed in the FC as $Z_{\text{spec}}$ (classical) or $\Phi$ (proposed). When classical verification succeeds, $Z_{\text{impl}}$ lies in the implementability overlap (BC $\cap$ FC via homomorphism $h$). Assertional checking may apply to buildable $Z'$ outside that overlap if $Z' \models_{\mathcal{L}} \Phi$.}
  \label{fig:t3sd-cotyledons}
\end{figure}

We restrict attention to the subset of system designs that can be expressed within Wymore's discrete systems (DSYSTEMS): a deterministic 5-tuple $Z = (S, I, O, \delta, \rho)$ with state set $S$, input set $I$, output set $O$, transition function $\delta$, and readout $\rho$ \parencite{Wymore1993}. Time index $t \in T$ advances in discrete, atomic ticks while state components may be continuous-valued at each tick (e.g., $x,y \in \mathbb{R}$).  ``Discrete system'' here refers to the clock, not to finite or discrete-valued states. Verification in this paper evaluates properties at tick boundaries, using the same convention as Wymore's readout function $\rho$. Wymore explicitly allows $S$, $I$, and $O$ to be \emph{not-necessarily-finite}, i.e., possibly infinite, compared to finite state machines only. Given an input trajectory $f: T \to I$ and initial state $s_0 \in S$, Wymore's trajectory theorems yield unique state and output trajectories $g: T \to S$ and $y: T \to O$. At ticks without an external command, we treat $f(t)$ as hold-last-input or as an environment-defined value in $I$ so that $\delta(g(t), f(t))$ is defined. Systems of interest here are open: verification is performed at the system boundary through observable inputs and outputs, consistent with Systems Theory \parencite{salado2019mbse}.

\miniheading{Notation}
Table~\ref{tab:notation} maps Wymore textbook notation \parencite{Wymore1993} to the functional notation used in this paper.

\begingroup
\renewcommand{\arraystretch}{1.4}
\footnotesize
\begin{tabularx}{\textwidth}{@{}A{0.20\textwidth} A{0.32\textwidth} >{\raggedright\arraybackslash}X@{}}
    \rowcolor{tableheader}
    \headingfont\bfseries Name & \headingfont\bfseries Wymore & \headingfont\bfseries Mathematical notation \\
    \addlinespace[3pt]
    Discrete system 5-tuple & $(SZ, IZ, OZ, \text{NZ}, \text{RZ})$ & $(S, I, O, \delta, \rho)$ \\
    Next state function & $\text{NZ}(x, p)$ & $\delta(s, i)$ \\
    Readout function & $\text{RZ}(x)$ & $\rho(s)$ \\
    Time scale & $\text{TZ}$ & $T$ \\
    Input trajectories & $\text{ITZ}$ & $f : T \to I$ \\
    Initial state & $x$ in $\text{STZ}(f, x)$ & $s_0$ \\
    State trajectory & $\text{STZ}(f, x)$ & $g : T \to S$ \\
    State at time $t$ & $\text{STZ}(f, x)(t)$ & $g(t)$ \\
    Output trajectory & $\text{OTZ}(f, x)$ & $y : T \to O$ \\
    Output at time $t$ & $\text{OTZ}(f, x)(t)$ & $y(t)$ \\
    System experiments & $\text{EXZ}$ & $(f, s_0, t)$ \\
    Homomorphism maps & $\text{HS}, \text{HI}, \text{HO}$ & $h_S, h_I, h_O$ \\
    Next-state consistency & $\text{HS}(\text{NZ1}(x,p)) = \text{NZ1}(\text{HS}(x), \text{HI}(p))$ & $h_S(\delta_2(s,i)) = \delta_1(h_S(s), h_I(i))$ \\
    Readout consistency & $\text{HO}(\text{RZ2}(x)) = \text{RZ1}(\text{HS}(x))$ & $h_O(\rho_2(s)) = \rho_1(h_S(s))$ \\
\end{tabularx}
\captionof{table}{Wymore-to-functional notation map.}
\label{tab:notation}
\endgroup
\vspace{4pt}

We use the following conventions throughout:
\begin{itemize}
    \item \textbf{DSYSTEM / discrete system:} deterministic 5-tuple $Z=(S,I,O,\delta,\rho)$; the time scale $T$ is discrete, but $S$, $I$, and $O$ may be infinite.
    \item \textbf{Reference system $Z_{\text{spec}}$:} a specification-system 5-tuple in the FC that expresses functional intent for constructive verification \parencite{Wymore1993} (synonymous with ``specification system'' in figures).
    \item \textbf{Wymore system homomorphism $h$:} surjective maps $h_S,h_I,h_O$ on states, inputs, and outputs that preserve $\delta$ and $\rho$ tick-wise; we write $h: Z_{\text{impl}} \to Z_{\text{spec}}$.
    \item \textbf{Property set $\Phi$:} temporal formulas in dialect $\mathcal{L}$ stating mission intent; \textbf{execution encoding $\psi_Z$:} FO-LTL formula from $\texttt{compileSystemFO}(Z,s_0)$ capturing dynamics of 5-tuple $Z$ (Section~\ref{sec:dsystem-encoding}).
    \item \textbf{Encode} maps natural-language requirements to $\Phi$; \textbf{formalize} maps them to $Z_{\text{spec}}$; \textbf{compile} is reserved for $\texttt{compileSystemFO}$ (5-tuple $\to \psi_Z$), not for spec$\to\Phi$ translation.
    \item \textbf{$Z \models_{\mathcal{L}} \Phi$:} FC membership---5-tuple $Z$ satisfies $\Phi$ under all admissible input trajectories. \textbf{Verification} of a candidate $Z_{\text{impl}}$ is the membership test $Z_{\text{impl}} \models_{\mathcal{L}} \Phi$.
    \item \textbf{Property checking} is evaluating $Z_{\text{impl}} \models_{\mathcal{L}} \Phi$. \textbf{Model checking} refers specifically to automated LTL verification on finite-state abstractions \parencite{BaierKatoen2008}.
\end{itemize}
The parameter $\mathcal{L}$ may range over temporal dialects; FO-LTL is used for $\Phi$, for $\psi_Z$, and in the encoding theorem. Propositional LTL is the finite-state tooling subset when $|S|$, $|I|$, $|O|$ are small.

\subsection{Temporal Logics}
Temporal logics are formal languages for stating requirements over time-indexed execution traces \parencite{BaierKatoen2008, Pnueli1977}. They are used to specify safety (forbidden conditions never occur), liveness (desired conditions eventually occur), and ordering constraints on behavior.

\begingroup
\renewcommand{\arraystretch}{1.35}
\footnotesize
\begin{tabularx}{\textwidth}{@{}P{0.10\textwidth} P{0.14\textwidth} A{0.30\textwidth} A{0.40\textwidth}@{}}
    \rowcolor{tableheader}
    \headingfont\bfseries Op. & \headingfont\bfseries Name & \headingfont\bfseries Example & \headingfont\bfseries Reading (over trace $\sigma$) \\
    \addlinespace[2pt]
    $G\,\varphi$ & always & $G\,\text{safe}$ & $\varphi$ holds at every tick \\
    $F\,\varphi$ & eventually & $F\,\text{done}$ & $\varphi$ holds at some future tick \\
    $X\,\varphi$ & next & $X\,\texttt{ack}$ & $\varphi$ holds at the next tick \\
    $\varphi \rightarrow \psi$ & implication & $\text{start} \rightarrow F\,\text{complete}$ & if $\varphi$ holds, then $\psi$ eventually holds \\
\end{tabularx}
\captionof{table}{Common LTL operators and illustrative readings.}
\label{tab:ltl-operators}
\endgroup
\vspace{4pt}

Table~\ref{tab:ltl-operators} summarizes some operators used in this paper. First-Order LTL (FO-LTL) adds quantifiers over state, input, or output variables at each tick and are needed when requirements relate continuous variables or indexed elements. Other dialects (e.g., Signal Temporal Logic (STL) for continuous signals) remain admissible as $\mathcal{L}$. Because functional intent varies in complexity from discrete event sequencing to continuous signal bounds, our framework parameterizes the chosen logic as $\mathcal{L}$ which allows arbitrarily complex property specification, which may create a larger possible FC than is generated by DSYSTEMS alone.

\subsection{Related Formalisms}
Generic state model checking does not, by itself, supply Wymore trajectory semantics or a constructive-versus-assertional account of T3SD verification. Hybrid and event-scheduled formalisms address dynamics outside Wymore's discrete systems with uniform, atomic clock \parencite{platzer2018logical,wach2021devs}.
